\cotefile{156-4}
\scanpage{74}\mspage{73}

\deleted{\uline{Corollaire} \unsure{2} \unsure{La}}

\uline{\unsure{Définition|Definition}} On dit \deleted{\unsure{alors}} \inserted{que $G$ \uline{raffine} $F$} \deleted{\unsure{si|dit} \unsure{est un} \unsure{raffint|raffinement}} ou est un \uline{raffinement} de $F$ \deleted{\unsure{.} On écrit \unsure{alors}} si $\deleted{\unsure{\in}}\exists$ un $F'$ avec \deleted{$G\ll F$} $G\leq F'\preceq F$. On écrit $G\ll F$

\uline{Corollaire} \deleted{\unsure{2}} \unsure{la|La} relation $G\ll F$ est une relation d'ordre (impliquée par $G\leq F$ et par $G\preceq F$)
\margin{\uline{\unsure{Utilise|Utiliser} $C6$,} \unsure{ni} \uline{$C5$ ni $C7$}.}

Transitivité : Soit $F\deleted{\unsure{\leq}}\ll G\ll H$, donc

\diagram{\begin{tikzcd}[ampersand replacement=\&,row sep=0.1em,column sep=0.3em]
 {} \& \deleted{G} \& {} \& \deleted{H'} \& \deleted{\preceq} \& \deleted{H} \\
 {} \& {} \& {} \& \deleted{\rotatebox{90}{$\leq$}} \& {} \& {} \\
 \deleted{H} \& \deleted{G'} \& \deleted{\preceq} \& \deleted{G} \& {} \& {} \\
 {} \& \deleted{\rotatebox{90}{$\leq$}} \& {} \& {} \& {} \& {} \\
 {} \& \deleted{\unsure{H|G}}\deleted{F} \& {} \& {} \& {} \& {}
 \arrow[from=1-2,to=3-1,dash]\arrow[from=1-4,to=4-2,dash]
 \arrow[from=1-5,to=3-3,dash]\arrow[from=1-6,to=3-5,dash]
\end{tikzcd}}

\diagram{\begin{tikzcd}[ampersand replacement=\&,row sep=0.08em,column sep=0.22em]
 H \& \succeq \& H' \& {} \& {} \\
 \vdots \& {} \& \rotatebox{90}{$\leq$} \& {} \& {} \\
 {} \& \cdots \& G \& \succeq \& G' \\
 {} \& {} \& {} \& {} \& \rotatebox{90}{$\leq$} \\
 {} \& {} \& {} \& {} \& \deleted{\unsure{H}}F
\end{tikzcd}}

on complète par $\unsure{6\text{b}|6\text{a}|C6|C5}$

\diagram{\begin{tikzcd}[ampersand replacement=\&,row sep=0.08em,column sep=0.22em]
 H \& \succeq \& H' \& {} \& {} \& {} \\
 \rotatebox{90}{$\leq$} \& \circ \& \rotatebox{90}{$\leq$} \& {} \& {} \& {} \\
 \bar H \& \succeq \& G \& \succeq \& G' \& \quad (*) \\
 \rotatebox{90}{$\leq$} \& \circ \& \rotatebox{90}{$\leq$} \& \circ \& \rotatebox{90}{$\leq$} \& {} \\
 K \& \succeq \& \bar G \& \succeq \& F \& {}
\end{tikzcd}}

et on gagne, \unsure{car|on|ou}
\[
 F\preceq K\leq H
\]

Supposons maintenant que $F=H$ i.e. $F\ll G$ et $G\ll F$, prouvons $F=\deleted{\unsure{H}}G$. On \deleted{\unsure{tire de}}
\[
 \begin{aligned}
 F\preceq K\leq{} & F & \text{on va en déduire}\\[-2pt]
 &\rotatebox{90}{$= $}\\[-2pt]
 &H & K=F
 \end{aligned}
\]
et comme $\underset{\substack{\rotatebox{90}{$=$}\\F}}{K}\leq\bar H\leq F$ \unsure{etc} on obtiendra
\[
 \bar H=F\quad \deleted{\text{et de m\textsuperscript{me}}\
 \unsure{\bar G|\bar H}=\unsure{F|K}\ (\text{comme } F\preceq\bar G\preceq
 \underset{\substack{\rotatebox{90}{$=$}\\F}}{K})\ \bar G=F}
\]
donc $G\preceq F$. De façon symétrique on aura $F\preceq G$, d'où $F=G$. Donc il reste à prouver le

\uline{Lemme} Soient $F,K\in\unsure{\mathcal F|\mathscr F}$ avec $F\preceq K\leq F$. Alors $K=F$ (i.e. si $F$ est une subdivision d'une sous-figure $K$, alors $F=K$)
\scanpage{75}\mspage{74}

Cela \unsure{résultera|résulte} du Corollaire qui suit, et de $C5$

\uline{Corollaire} (Propriété de la relation $\leq$ !) Soient $F$ une figure, $K$ une sous-figure \deleted{qui} telle que $K\neq F$. Alors il existe une \deleted{\unsure{str}} figure $G$ (une \unsure{multistrate}, si on veut) qui est \deleted{\unsure{disj}} disjointe de $K$, mais non de $F$. En d'autres \unsure{termes}, $\operatorname{supp} K\neq\operatorname{supp} F$.

Ce n'est pas ``formel'', en termes des axiomes $C1$ \unsure{:|;} $C4$ sur $(\unsure{\mathcal F|\mathscr F},\leq)$, supposons p.ex. que $\mathcal F$ admette un plus grand élément, donc il est défini comme l'ens \unsure{de|des} \uline{toutes} les parties fermées de l'ens ordonné $\unsure{M|\mathcal M}$, et ``disjoint'' signifie ici $A\cap B=\varnothing$ sans plus (la relation de compatibilité $R$ est la relation \unsure{tjrs|tjrs.} vraie). Mais si $F$ est une partie fermée, $K\subset F$ une sous partie fermée \inserted{$\neq F$}, il n'est pas clair qu'on puisse trouver une partie fermée $G$ telle que $G\cap\deleted{F}K=\varnothing$, $G\cap F\neq\varnothing$. \deleted{P.ex. $M=[0,\unsure{N|\mathbb N}]$, \unsure{$F$} avec \unsure{rel\textsuperscript{n}} d'ordre \unsure{(?)|ord} habituel, $F=M$, $K=[0,1]$,} \inserted{\unsure{Si|En} p.ex. $M$ a un plus petit élément $0$,} Alors pour \uline{toute} partie fermée non vide, les $G$ tels que $F$ et $G$ disjoints sont les $G$ vides, donc ça ne change pas si on remplace $F$ par \deleted{\unsure{F}}$K\unsure{\in|\subset}F$, \unsure{tant|\deleted{tant}} que $K$ n'est pas vide
\scanpage{76}\mspage{75}

Dans ce cas, l'axiome $C6$ (pour une relation \inserted{d'ordre} \unsure{suppl|suppl.} $\preceq$ sur $\unsure{\mathcal F|\mathscr F}=\unsure{\mathrm P|\mathcal P}_{f}(\unsure{M|\mathcal M})$) signifie \unsure{seulement|seulmt} que si $F'\preceq F$, $F'$ est vide ssi $F$ l'est. Prenons p.ex. \deleted{$\unsure{F|K}\unsure{\preceq}$} $\preceq$ \deleted{\unsure{égale}} \unsure{à|:} la relation d'ordre discrète, c'est OK. \deleted{$C7$ \unsure{a} est \unsure{tautologique}} $C6,C7$ tautologiques. \unsure{On voit donc|On veut dire} qu'il manque un axiome, pour assurer la validité du corollaire. On devine que ce sera la partie \unsure{a)|--)} du \unsure{Const|Cons} $4$ (p. 35), \unsure{jointe} à Const $7$. Donc pour l'instant je ne vais pas utiliser que le \unsure{préordre} $\ll$ soit un ordre.
% slanted
\margin{Aussi, cela résulte \unsure{aussitôt|aussi} \unsure{de la} \unsure{modification} de $C4$ \unsure{ancien} (page $47$).}

\deleted{\unsure{Je} songe à la propriété \unsure{suiv\textsuperscript{te}|suiv.}, qu'il faudrait avoir tôt ou tard :}

\deleted{\deleted{Soient $\unsure{Y|y},\unsure{X|x}\in M$ \unsure{avec} $Y\ll X$, et} supposons que \deleted{$Y'\leq Y$, $Y'\in\partial X$, $Y\cap\partial X=$} \deleted{$\unsure{\exists|\nexists}\,Y'\leq Y$, \unsure{alors}} $Y'\leq Y$ $(Y'\in M)$ tels que \deleted{$\unsure{\exists}\,Y'\ll\partial X$}}

\deleted{Soient $F$ figure, $F'$ sous-figure, $G\ll F$, $G'$ \inserted{$\leq G$} le \unsure{raffint|raffinement} de $F'$ induit par $G$ (cf plus bas). \deleted{supposons que} \inserted{Alors} $G'=\varnothing$. \deleted{Alors $G$ est} ssi $G$ disjoint de $F'$.}

Il faut développer la notion de \uline{\unsure{raffinement|raffinements} induit}. J'ai envie de \deleted{\unsure{démontrer}} \inserted{\unsure{prouver}} l'énoncé (quitte à rajouter des axiomes) :
\scanpage{77}\mspage{76}

\deleted{\uline{\unsure{Prop.|Déf.}}} \inserted{\uline{Prop.}} Soit

\diagram{\begin{tikzcd}[ampersand replacement=\&,row sep=0.15em,column sep=0.3em]
 G \& \ll \& F \\
 \deleted{\rotatebox{90}{$\leq$}} \& {} \& \rotatebox{90}{$\leq$} \\
 \deleted{\unsure{G'|G?}} \& {} \& F'\,\unsure{.}
\end{tikzcd}}

alors il existe \unsure{un|une} plus grande sous figure $F'$ de $G$, qui raffine $F'$ (et qu'on appelle raffinement de $F'$ induit \deleted{par le} sur la sous figure $F'$ par le raffinement $G$ de $F$, ou la sous-figure du \deleted{$G$ \unsure{induit}} \unsure{raff\textsuperscript{t}|raff} $G$ de $F$, induite par la sous-figure $F'$ de $F$). \deleted{\unsure{Elle est}} \unsure{Ses|Les} strates \deleted{(c'est $\unsure{\pm|=}$ tautologique)} sont les strates de $G$ qui raffinent $F'$. \deleted{De plus, $G$ et $F$}

\diagram{\begin{tikzcd}[ampersand replacement=\&,row sep=0.12em,column sep=0.3em]
 G \& \ll \& F \\
 \rotatebox{90}{$\leq$} \& {} \& \rotatebox{90}{$\leq$} \\
 G' \& \ll \& F'
\end{tikzcd}}

De plus, $G'$ est un $\operatorname{Inf}$ de $G$ et $F'$ pour la relation $\underset{\unsure{\cdot}}{\ll}$, en d'autres termes, si \deleted{\unsure{$X$} $G$ est} $K$ est une figure telle que $K\ll G$ et $K\ll F'$, alors $K\ll G'$.

\margin{Dire (ou poser) que $G'=\varnothing\Longleftrightarrow G$ disjoint de $F'$ (seul $\Rightarrow$ non évident\unsure{....|...})}

\deleted{\uline{\unsure{Dém|Démon}} Soit $H$ une sous-figure \deleted{de $G$ telle que} \inserted{de $H$} $H\ll F'$, alors pour tte strate $X\leq H$, on aura $X\ll F'$, donc \unsure{Si} on désigne par $G'$ la sous-figure de $G$ engendrée par des strates $Y\leq G$ telles que $Y\ll H$, alors \deleted{$\unsure{X}$} $H\leq G$ puisque}
\[
 \deleted{\displaystyle H=\mathop{\mathrm{Sup}}\limits_{\substack{X\in\deleted{\unsure{H|G}}\,\unsure{M|\mathcal M}\\X\leq H}}X\deleted{\unsure{<}}\leq\mathop{\mathrm{Sup}}\limits_{\substack{X\in M\\X\in G'}}X=G'\unsure{.|,}}
\]
\deleted{D'autre part, $G'\ll$}
\scanpage{78}\mspage{77}

On a par \unsure{définition|definition} de $\ll$
\[
 G\leq\bar F\preceq F
\]
\margin{$G$ sous-figure d'une subd. $\bar F$ de $F$}
et on va appliquer $C6$

\diagram{\begin{tikzcd}[ampersand replacement=\&,row sep=0.07em,column sep=0.25em]
 G \& \leq \& \bar F \& \preceq \& F\\
 \rotatebox{90}{$\leq$} \& {} \& \rotatebox{90}{$\leq$} \& {} \& \rotatebox{90}{$\leq$}\\
 G' \& \leq \& \bar F' \& \preceq \& F'\\
 \rotatebox{90}{$=$} \& {} \& {} \& {} \& {}\\
 G\cap\bar F' \& {} \& {} \& {} \& {}
\end{tikzcd}}

\deleted{\unsure{Soit}} Prouvons que $G'$ est un $\inf$ pour $\ll$

Soit $K$ telle que $K\ll G$, $K\ll F'$ on doit prouver $K\ll G'$. Pour ceci, on prouvera d'abord $K\ll\bar F'$, et on est ramené aux deux cas particuliers où on a soit $G\preceq F$, soit $G\leq F$. Donc

\uline{Lemme} 1 Soit

\diagram{\begin{tikzcd}[ampersand replacement=\&,row sep=0.08em,column sep=0.25em]
 G \& \preceq \& F \& \quad\text{(cf } C6\text{)}\\
 \rotatebox{90}{$\leq$} \& {} \& \rotatebox{90}{$\leq$} \& {}\\
 G' \& \preceq \& F' \& {}
\end{tikzcd}}

et soit $K$ \unsure{tel|tels|t.q.} que $K\ll G$ et $K\ll\deleted{\unsure{G}}F'$, alors $K\ll G'$.

En effet, par \deleted{\unsure{le}} \inserted{déf.} on aura $K\leq\bar G\preceq\deleted{\unsure{F}}G$

\diagram{\begin{tikzcd}[ampersand replacement=\&,row sep=0.08em,column sep=0.25em]
 K \& \leq \& \bar G \& \preceq \& \deleted{F}G \& \preceq \& F\\
 {} \& {} \& \rotatebox{90}{$\leq$} \& {} \& \rotatebox{90}{$\leq$} \& {} \& \rotatebox{90}{$\leq$}\\
 {} \& {} \& \bar G' \& \preceq \& G' \& \preceq \& F'
\end{tikzcd}}

pour prouver $K\ll G'$, il suffit de prouver $K\leq\overline{G'}$. \deleted{\unsure{Or}} \deleted{Soit $X\in\unsure{M|\mathcal M}$, $X\in\unsure{G|C}$} \inserted{ou encore, pour} $X\in M$,
\[
 \underset{\inserted{\text{d'où } X\leq\bar G,\ X\ll F'}}{\underline{\rule[-.5ex]{.3pt}{.45ex}\,X\leq K\,\rule[-.5ex]{.3pt}{.45ex}}}\ \Rightarrow X\leq\bar G',
 \quad\text{\unsure{or|ou}}
\]
Or cela résulte de $C6$, $(b)$ appliqué à

\diagram{\begin{tikzcd}[ampersand replacement=\&,row sep=0.08em,column sep=0.25em]
 \bar G \& \preceq \& F\\
 \rotatebox{90}{$\leq$} \& {} \& \rotatebox{90}{$\leq$}\\
 \bar G' \& \preceq \& F'
\end{tikzcd}}
\scanpage{79}
\mspage{78}

\uline{Lemme 2} Soit
\diagram{\begin{tikzcd}[ampersand replacement=\&,row sep=0.06em,column sep=0.20em]
G \& \leq \& F\\
\rotatebox{90}{$\leq$} \& \& \rotatebox{90}{$\leq$}\\
G' \& \leq \& F'\\
\rotatebox{90}{$=$} \& \& \\
G\cap F' \& \&
\end{tikzcd}}
et soit $K$ tel que $K\ll G$ et $K\ll F'$, alors $K\ll G'$.

On écrit encore $K\leq\bar G\preceq G$, et \ldots

\deleted{\unsure{on}} J'abandonne ! Il vaut mieux se reporter au texte page \unsure{51|61}.

\margin{\uuline{NB} $F'\unsure{\mathrel{\vert\!-\!\vert}|\mapsto}G'$ \quad $\text{Sous fig }F\unsure{-|\to}\text{ Sous fig }G$ commute aux $\operatorname{Sup}_{\leq}$ \unsure{quelc.}, \deleted{\unsure{puis}} et aux \deleted{\illegible{1}} $\operatorname{Inf}_{\leq}$ quelconques\ldots}

\deleted{Soit $F$ une figure}

\deleted{\uline{\unsure{C8|Co\,8}} Soit $F$ une figure, et $X$ une multistrate telle que $X\ll F$ (i.e $X\leq F'\preceq F$). Alors $\exists\,Y\,\unsure{\in|i}$, \deleted{\unsure{multi}}\unsure{strate} de $F$, telle que $X\deleted{\unsure{\leq}}\ll Y$.}

\uline{Proposition} Soient $F$ une figure, $X\in\unsure{M|\mathcal M}$, tels que
\[
 X\ll F.
\]

a) \deleted{Alors, \unsure{je} $\exists\,Y\in$} $\exists\,Y\in M$, $Y\leq F$ telle que
\[
 X\deleted{\unsure{\leq}}\ll Y
\]

b) Parmi les $Y$ satisfaisant \unsure{cette|à cette} condition, il y un plus petit élément, \deleted{\unsure{donné}} \deleted{\unsure{il est}} \inserted{\unsure{J|$\hookleftarrow$} \unsure{C avec|avec|car} $\unsure{C12|C\,12}$.} \deleted{caractérisé par} caractérisé par la propriété que la sous-figure de $X$ induite par $\partial Y$ est $\neq X$.

\uline{Démonstration} a) $F=\mathop{\mathrm{Sup}}\limits_{Y_i\in\bar F}\deleted{\unsure{X_i|Y_i}}Y_{i\deleted{\unsure{\leq|i}}}$, donc \deleted{\unsure{$X$}} (par $C6$ $(c)$) $X=\mathop{\mathrm{Sup}}\limits_{i}G_i$ où $G_i$ est la sous fig. de $X$ induite par les $\deleted{\unsure{X_i}}Y_i$, donc $\exists\,i$ tel que $X\in G_i$, donc $X\ll Y_{\unsure{i|v}}$.
\scanpage{80}
\mspage{79}

Parmi les $Y$\inserted{$\in\bar F$} tels que $X\ll Y$, choisissons un élément minimal (cela existe, si on admet que la dim. combinatoire des figures \unsure{irréd|irréd.} (\unsure{i.e.|$=$.} multistrates) est finie -- on y reviendra sous forme d'un \unsure{énoncé} plus fort encore \ldots\ \deleted{\unsure{Donc} \unsure{Donc !}}

Dire que $Y$ est minimal signifie que pour $Y'$ \deleted{\unsure{multi}}strate de $Y$, on n'a pas $X\deleted{\unsure{\leq}}\ll Y'$ i.e. $X\neq F_{Y'}$, sous-figure de $X$ induite par $Y'$. Mais cela signifie aussi que l'image inverse de $\partial Y=\mathop{\mathrm{Sup}}\limits_{\substack{\deleted{\unsure{Y'\leq\partial Y|Y'\leq Y}}\\Y'\in\unsure{M|\mathcal M}\\Y'\leq Y}}Y'$ est \unsure{différent|différente} \unsure{de|\deleted{de}} $X$ (\unsure{appliquant|appliquer} à nouveau $C6$ $(c)$). \deleted{Si on a}

Si \unsure{on} a $Y'\in\bar F$ tel que $X\ll Y'$, prouvons que $\deleted{\unsure{X|\leq}}Y\leq Y'$. \deleted{\unsure{On}} En \deleted{\unsure{appliquant|appliquer} le} effet l'image inverse $F_{Y'}=X$, \unsure{or}
\[
 \underset{\inserted{\substack{\rotatebox{90}{$=$}\\X}}}{F_Y}\cap
 \underset{\inserted{\substack{\rotatebox{90}{$=$}\\X}}}{F_{Y'}}=F_{Y\cap Y'}
\]
à cause des propriétés formelles du $\operatorname{Inf}$ pour $\ll$ donc $F_{Y\cap Y'}\deleted{\unsure{F_Y\cap F_{Y'}|F_{Y'}}}=X$, i.e. $X$ se factorise par $Y\cap Y'$ donc par un $Z\leq Y\cap Y'\leq Y$ donc $Z=Y$ (\unsure{condition|caractère|car. être|car. étant} \unsure{minimal|minimale}) et par suite $Y=Y'$, on gagne !!
\scanpage{81}
\mspage{80}

\uline{Corollaire} Soient $F,G$ avec $F\ll G$. Pour tout $X\in\bar F$, soit $\unsure{\varphi|\psi|\phi}(X)\in\bar G$ la plus petite \deleted{multi}strate \inserted{$Y$} de $G$ telle que $\deleted{\varphi(X)\ll}\,X\ll Y$. L'application
\[
 \varphi:\bar F\longrightarrow\bar G
\]
ainsi obtenue est \uline{croissante} (trivial).

Jusqu'à \unsure{maintnt|maintenant}, les axiomes \deleted{\unsure{pour}} $C5$, $C6$, $C7$ introduits pour $\preceq$ sont tautologiques, si on prend pour $\preceq$ l'égalité, \inserted{alors $\ll$ est $\leq$, \unsure{i.e.|scil.} l'inclusion de figures.} Il serait temps d'en venir à des propriétés \deleted{telles que} plus spécifiques !

\uline{14 juin}

\margin{Je vais poser $\widetilde F$ \unsure{ald|a l d} $\bar F$ pour l'ens des $X\in\unsure{M|\mathcal M}$ tels que $X\leq F$ (pour éviter confusion avec notion d'\unsure{adhérence|adhésion})}

\deleted{\uline{Proposition}} \inserted{Corollaire de prop. \unsure{pag. 76|page 76}} (cf l'\uline{ancien} $C4$ \deleted{\unsure{pa}} page 47)

Soient $G'\deleted{\unsure{\leq}}\unsure{\ll|\leq}G\deleted{\unsure{\leq}}\unsure{\preceq|\leq}F$ alors $\widetilde{G'}\cap\widetilde F\subset\widetilde G$
\inserted{(donc $\widetilde{G'}\cap\widetilde F=\widetilde{G'}\cap\widetilde G$)}

En d'autres termes, toute \inserted{(multi)}strate \deleted{\unsure{de}} \deleted{\unsure{figure}} \inserted{$X$ qui raffine la sous-figure $G$ de $F$ \unsure{et}} est \inserted{(multi)}strate de $F$ est \inserted{(multi)}strate de $G$.

\deleted{\deleted{Soit $X$ strate de $G'$ qui est strate de $F$}}

\deleted{\deleted{\unsure{D}} \unsure{Dém.} $X$ est une strate d'un raffinement $\bar G$ de $G$, \inserted{lequel est} \unsure{induit} par un raffinement $\bar F$ de $F$ ($C6$)}
\diagram{\begin{tikzcd}[ampersand replacement=\&,row sep=0.05em,column sep=0.16em]
\deleted{G} \& \deleted{\unsure{=|\leq}} \& \deleted{F} \& \deleted{\geq} \& \deleted{X}\\
\deleted{\rotatebox{90}{$\preceq$}} \& \& \deleted{\rotatebox{90}{$\preceq$}} \& \& \deleted{\unsure{;}}\\
\deleted{\bar G} \& \deleted{\leq} \& \deleted{\bar F} \& \& \\
\deleted{\rotatebox{90}{$\leq$}} \& \& \& \& \\
\deleted{X} \& \& \& \&
\arrow[from=1-2,to=5-1,dash]
\end{tikzcd}}
% slanted
\margin{\unsure{hi hi hi|si si si|hé hé hé} $\longrightarrow$}
\deleted{Supposons $X\leq F$, \deleted{et considérons les images inverses} \inserted{l'image inverse de}}
